% The Brexit negotiation as a game: tasks (i)-(vii) of Cases/Brexit/intrucoes.md.
% Compile: pdflatex brexit-game.tex (twice). Check: pdffonts (all Type 1, uni = yes).
% Every number is asserted by brexit-check.py in this folder. The case PDFs are not reproduced:
% only the cost figures listed in .agent/plans/cases-lecture-03-brexit.md, section 2, are used.
\documentclass[11pt,a4paper]{article}
\usepackage[utf8]{inputenc}
\usepackage[T1]{fontenc}
\usepackage{lmodern}   % Type 1 vector fonts (replaces PK bitmaps)
\usepackage{cmap}      % ToUnicode CMap -> copyable text and accents
\usepackage[portuguese,english]{babel}
\usepackage{amsmath,amssymb,amsthm,mathtools}
\usepackage[a4paper,margin=2.4cm]{geometry}
\usepackage{microtype}
\usepackage{booktabs,array,tabularx}
\usepackage{enumitem}
\usepackage{xcolor}
\usepackage{tikz}
\usepackage[most]{tcolorbox}
\usepackage{fancyhdr}
\usepackage{caption}
\usepackage[hidelinks]{hyperref}

\setlist{topsep=3pt,itemsep=2pt}
\setlength{\parindent}{0pt}
\setlength{\parskip}{5pt}
\renewcommand{\arraystretch}{1.2}
\allowdisplaybreaks
\captionsetup{font=small,labelfont=bf}

\pagestyle{fancy}\fancyhf{}
\lhead{\footnotesize Economics of Organizations and Incentives --- MPE Insper}
\rhead{\footnotesize Case: the Brexit negotiation as a game}
\cfoot{\footnotesize\thepage}
\renewcommand{\headrulewidth}{0.3pt}
\setlength{\headheight}{14pt}

\definecolor{ans}{RGB}{20,110,70}
\definecolor{src}{RGB}{30,80,150}
\definecolor{warn}{RGB}{190,90,20}
\newcommand{\answer}[1]{\begin{tcolorbox}[colback=ans!5,colframe=ans,boxrule=0.6pt,
  left=5pt,right=5pt,top=3pt,bottom=3pt]\textbf{Answer.} #1\end{tcolorbox}}
\newcommand{\gbp}{\pounds}
\newcommand{\Nt}{\mathit{Nt}}
\newcommand{\dF}{\mathit{F}}
\newcommand{\dU}{\mathit{U}}

% One UK decision node with its Present/Reject edges, the vote and three leaves.
% #1 index of its top leaf, #2 parent EU node, #3 deal label, #4-#6 payoffs (pass, fail, reject),
% #7 style of the EU edge, #8 style of the present edge, #9 style of the reject edge.
\newcommand{\branch}[9]{%
  \node[dec,label={[font=\tiny,inner sep=1pt]above:UK}] (u#1) at (4,{#1+1}) {};
  \node[nat,label={[font=\tiny,inner sep=1pt]above:Parl.}] (n#1) at (6,{#1+0.5}) {};
  \draw[#7] (#2) -- node[lab]{#3} (u#1);
  \draw[#8] (u#1) -- node[lab,above,sloped]{present} (n#1);
  \draw[#9] (u#1) -- node[lab,below,sloped,pos=0.4]{reject} (8,{#1+2});
  \draw[e] (n#1) -- node[lab,above,pos=0.6]{$0.8$} (8,#1);
  \draw[e] (n#1) -- node[lab,below,pos=0.6]{$0.2$} (8,{#1+1});
  \node[right] at (8,#1) {#4};
  \node[right] at (8,{#1+1}) {#5};
  \node[right] at (8,{#1+2}) {#6};
}

\title{\textbf{Boris's Brexit Dilemma as a Game}\\[4pt]
  \large Players, tree, subgame-perfect equilibrium and the value of commitment\\[2pt]
  \normalsize Case write-up for Lecture 3 --- Economics of Organizations and Incentives (MPE Insper)}
\date{November 2026}

\begin{document}
\maketitle
\thispagestyle{empty}

\begin{tcolorbox}[colback=src!4,colframe=src,boxrule=0.5pt,title={The assignment and the sources}]
\textbf{Cases.} ``Boris's Brexit Dilemma (A)'' and ``Boris's Brexit Dilemma (B)'' (Darden,
UV9621 and UV9624). The case text is not reproduced here. The model uses only the cost figures
summarized in \S4; every other element (the political cost $r$, the deterrence premium $k$, the
variants) is an addition of this write-up and is labelled as such.

\textbf{Assignment} (\texttt{Cases/Brexit/intrucoes.md}, quoted in the original Portuguese). In
pairs, students negotiate a Brexit deal with their counterpart and write the negotiation as a
game; in class each side decides whether to prepare for a no-deal Brexit (UK), which deal to offer
(EU) and whether to take the EU's proposal to Parliament (UK). The task list:

\begin{quote}\small\foreignlanguage{portuguese}{%
Para modelar a negociação como um jogo, seu grupo deve:
(i) identificar os jogadores;
(ii) indicar a ordem das jogadas e as ações disponíveis para cada jogador em cada ponto de
decisão;
(iii) indicar o que cada jogador sabe quando joga (por exemplo, a UE observa se o Reino Unido se
preparou para um Brexit sem acordo?);
(iv) atribuir payoffs a cada resultado final com base nas informações do caso (os números exatos
importam menos do que uma ordenação que vocês consigam justificar);
(v) desenhar a árvore do jogo; (vi) resolver o jogo por indução retroativa e encontrar o
equilíbrio perfeito em subjogos; e
(vii) comparar o equilíbrio com o resultado da negociação de vocês e, se forem diferentes,
explicar por quê.}
\end{quote}

\textbf{Method.} Extensive form, strategies, backward induction and subgame perfection: Tadelis
(2013) Ch.~7 (pp.~129--150) and Ch.~8 (pp.~151--174); credibility and commitment: Cabral (2017)
\S7.2, as in \texttt{Notas/L03/L03-06-credibility-entry-commitment.md}. Every number below is
asserted with exact fractions by \texttt{Cases/Brexit/brexit-check.py}, which also runs backward
induction programmatically on the tree.
\end{tcolorbox}

\begin{tcolorbox}[colback=ans!5,colframe=ans,boxrule=0.8pt,title={Executive summary}]
\textbf{Subgame-perfect equilibrium} (baseline: the EU sees whether the UK prepared; $r=0$):
\begin{itemize}[leftmargin=1.4em]
  \item \textbf{UK:} \emph{prepare} for no deal; if prepared, present only the favourable deal
    and reject the neutral and unfavourable ones; if unprepared, present every deal.
  \item \textbf{EU:} offer the \emph{favourable} deal if the UK prepared and the
    \emph{unfavourable} deal if it did not.
\end{itemize}
On the equilibrium path the UK prepares, the EU offers the favourable deal and the UK presents
it. Expected payoffs are $(-24.7,\,-14)$ in \gbp bn, against $(-28.6,\,-7.6)$ if the UK had not
prepared.

\textbf{Lesson.} A sunk, visible investment makes the threat to reject credible and shifts the
EU's offer. Preparation lowers the UK's ex-post cost of walking away from 35 to 22.5, below the
neutral and unfavourable deals. The EU, seeing this, concedes the favourable deal. The UK gains
3.9; the EU loses 6.4; the 2.5 difference is a deadweight loss.

\textbf{Robustness.} The result needs the preparation to be \emph{observed}: if it is hidden,
there is no pure equilibrium, and the UK prepares with probability $8/13$ while the EU offers the
favourable deal with probability $11/36$. Talks then break down with probability $50/117\approx0.43$.
It also needs the political cost of walking away to be small ($r<0.4$) and the EU's
deterrence premium on a favourable deal to be moderate ($k\le5$).
\end{tcolorbox}

\paragraph{Notation.}
\begin{itemize}[leftmargin=1.4em,topsep=1pt,itemsep=0pt]\small
  \item $d\in\{\dF,\Nt,\dU\}$: the soft-Brexit deal the EU offers (favourable, neutral,
    unfavourable to the UK); $c_d$ and $e_d$ are the UK's and the EU's costs of deal $d$.
  \item $H_{UK}=35$, $H_{EU}=18$: the costs of a hard (no-deal) Brexit.
  \item $\kappa=5$: the UK's preparation cost; $m=12.5$: the hard-Brexit cost it saves.
  \item $p=0.8$: the probability that Parliament passes a presented deal.
  \item $r\ge0$: the UK government's political cost of rejecting the EU's offer outright (not in
    the case's numbers). $k\ge0$: the EU's deterrence premium on a favourable deal (not in the case's
    numbers).
  \item Payoff $=-$cost, in \gbp bn. Players are risk neutral: they rank lotteries by expected
    cost.
\end{itemize}

%-------------------------------------------------------------------------------------------
\section{(i) The players}

\begin{itemize}
  \item \textbf{The United Kingdom}, represented by the government of Boris Johnson. It moves
    twice: it chooses whether to prepare for no deal, and whether to take the EU's offer to
    Parliament.
  \item \textbf{The European Union}, represented by its chief negotiator, Michel Barnier. It chooses
    which deal to offer.
  \item \textbf{Parliament, modelled as Nature.} A presented deal passes with probability
    $p=0.8$ and fails with probability $0.2$; a failed vote means a hard Brexit.
\end{itemize}

\textbf{Why Parliament is Nature.} The case data give the vote as a probability, not as the
choice of an agent with stated payoffs. The government cannot control it, and the EU cannot
bargain with it directly. A chance move with known probabilities is the standard way to put such
a step into an extensive form (Tadelis \S7.1, Definition~7.1, the ``Nature'' player). The model
also makes $p$ independent of the deal type, which is what a Nature move with one probability
implies.

\textbf{Remark: Parliament as a strategic player.} A strategic Parliament (say, a pivotal MP who
compares each deal with no deal) would make three changes. (1)~The game has three players, and
the vote becomes a best response, so $p$ depends on the deal: an unfavourable deal is more likely
to fail. \S7.4 shows that this \emph{helps} the UK, because the EU then has a reason to offer
more even to an unprepared UK. (2)~The government may want a \emph{tougher} Parliament: a veto
player that rejects bad deals is a commitment device. This is Schelling's (1960) point that the
power to bind oneself is bargaining power, and Putnam's (1988) ``two-level game''. (3)~With
several parties in Parliament, the vote is itself a bargaining game; the legislative bargaining of
\texttt{Notas/L04/L04-08-legislative-bargaining.md} would be the tool. We keep Parliament as
Nature in the baseline because it keeps the tree solvable by hand and isolates the commitment
effect.

%-------------------------------------------------------------------------------------------
\section{(ii) Timing and actions}

\begin{enumerate}[label=\textbf{Stage \arabic*.},leftmargin=4.6em]
  \item \textbf{UK} chooses \emph{Prepare} ($P$) or \emph{Not prepare} ($N$) for a no-deal
    Brexit. Preparation costs $\kappa=5$, is all-or-nothing and is sunk: it is paid even if a deal
    is later implemented.
  \item \textbf{EU} offers a soft-Brexit deal $d\in\{\dF,\Nt,\dU\}$.
  \item \textbf{UK} either \emph{presents} the deal to Parliament or \emph{rejects} it outright.
    Rejection means a hard Brexit, plus the political cost $r$.
  \item \textbf{Nature (Parliament)} passes a presented deal with probability $0.8$ (the deal is
    implemented) or rejects it with probability $0.2$ (hard Brexit, no political cost $r$:
    the government tried).
\end{enumerate}

\textbf{Strategies.} A pure strategy names an action at every decision node, including nodes that
are never reached (Tadelis Definition~7.4, pp.~138--139). The UK has one root node and six
present/reject nodes, one for each (preparation, offer) pair. The EU has two nodes, one for each
preparation choice. Counting by Tadelis's rule $\lvert S_i\rvert=\prod_k m_k$:
\begin{align*}
  \lvert S_{UK}\rvert &= 2\times2^6 = 128, & \lvert S_{EU}\rvert &= 3^2 = 9 .
\end{align*}
A UK strategy therefore reads like ``$P$; after $P$: present $\dF$, reject $\Nt$, reject $\dU$;
after $N$: present all three''. The EU's reads like ``$\dF$ after $P$, $\dU$ after $N$''.

%-------------------------------------------------------------------------------------------
\section{(iii) Information: who knows what}

\textbf{Baseline: the EU observes preparation.} Stockpiling, contingency contracts and border
infrastructure are visible, so the EU knows whether the UK prepared when it makes its offer.
The UK sees the offer before deciding whether to present it. Every information set is a
singleton. Because Nature moves, Tadelis's Definition~7.3 (p.~136) formally classifies the game as
one of imperfect information. Still, no player is ever unsure about an earlier move: Nature moves
last, and its probability is common knowledge. The UK presents without knowing how the vote will
go, but that is risk, not private information. We also take $r$ and $k$ as common knowledge in the
baseline; \S11 discusses what changes when they are not.

\textbf{Variant: the EU does not observe preparation.} The EU's two nodes (after $P$ and after
$N$) lie in one information set, drawn as a dashed line in Figure~\ref{fig:unobs}. The EU must choose
one offer for both. The UK still knows its own preparation and sees the offer, so its six
present/reject nodes remain singletons (perfect recall, Tadelis Definition~7.7, p.~142). The EU's strategy
set shrinks to $\{\dF,\Nt,\dU\}$. Strategically, stage~1 and stage~2 become simultaneous: the
same tree with a different information structure, as in the sequential and simultaneous Battle of
the Sexes of \texttt{Notas/L03/L03-01-extensive-form.md}~\S4. \S8 solves this variant.

%-------------------------------------------------------------------------------------------
\section{(iv) Payoffs}

\subsection{Inputs}

\begin{center}
\begin{tabular}{lcc}
\toprule
Outcome (costs in \gbp bn) & UK & EU \\
\midrule
Hard Brexit, UK unprepared & 35 & 18 \\
Soft Brexit, favourable deal $\dF$ & 19 & 13 \\
Soft Brexit, neutral deal $\Nt$ & 23 & 9 \\
Soft Brexit, unfavourable deal $\dU$ & 27 & 5 \\
\midrule
Preparation: cost $\kappa$, hard-Brexit cost saved $m$ & 5; \ 12.5 & --- \\
Probability Parliament passes a presented deal & \multicolumn{2}{c}{$p=0.8$} \\
\bottomrule
\end{tabular}
\end{center}

Preparation saves 2.5 per pound spent ($12.5=2.5\times5$), but only in a hard Brexit. In a soft
Brexit the 5 is lost. The EU's costs do not depend on whether the UK prepared.

\subsection{How each terminal payoff is computed}

A prepared UK's hard-Brexit cost is the unprepared cost, minus the saving, plus the sunk outlay:
\begin{align*}
  H_{UK}-m+\kappa &= 35-12.5+5\\
                  &= 22.5+5\\
                  &= 27.5 .
\end{align*}
A prepared UK in a soft Brexit pays the deal plus the outlay, $c_d+5$. A rejection adds $r$ to the
UK's hard-Brexit cost. The eighteen terminal nodes therefore carry only eight distinct payoff
pairs (Table~\ref{tab:leaves}).

\begin{table}[h]
\centering
\caption{Terminal payoffs (UK, EU) in \gbp bn. ``Fail'' is the vote failing; ``reject'' is the
UK rejecting outright.}\label{tab:leaves}
\begin{tabular}{llccc}
\toprule
& & \multicolumn{3}{c}{Offer $d$} \\
\cmidrule(lr){3-5}
UK prepared? & Event & $\dF$ & $\Nt$ & $\dU$ \\
\midrule
$N$ & deal passes   & $(-19,-13)$ & $(-23,-9)$ & $(-27,-5)$ \\
$N$ & vote fails    & $(-35,-18)$ & $(-35,-18)$ & $(-35,-18)$ \\
$N$ & UK rejects    & $(-35-r,-18)$ & $(-35-r,-18)$ & $(-35-r,-18)$ \\
\midrule
$P$ & deal passes   & $(-24,-13)$ & $(-28,-9)$ & $(-32,-5)$ \\
$P$ & vote fails    & $(-27.5,-18)$ & $(-27.5,-18)$ & $(-27.5,-18)$ \\
$P$ & UK rejects    & $(-27.5-r,-18)$ & $(-27.5-r,-18)$ & $(-27.5-r,-18)$ \\
\bottomrule
\end{tabular}
\end{table}

\subsection{Why the ordering is defensible}

The assignment stresses that a justified ordering matters more than exact numbers. Four ordinal
facts drive the whole solution; any numbers that preserve them give the same equilibrium.

\begin{enumerate}[label=(O\arabic*)]
  \item \textbf{An unprepared UK prefers every deal to no deal:} $27<35$. Even the unfavourable
    deal avoids the worst disruption.
  \item \textbf{A prepared UK prefers no deal to the neutral and unfavourable deals, but not to the
    favourable one:} $19<22.5<23<27$. Here 22.5 is the prepared hard-Brexit cost net of the sunk
    5, which is paid in either case.
  \item \textbf{The EU prefers any implemented deal to no deal, and less generous deals to more
    generous ones:} $5<9<13<18$. A hard Brexit hurts the EU too, but less than the UK
    ($18<35$): the EU is the larger party and the one with the better outside option.
  \item \textbf{Preparation plus the favourable deal beats no preparation plus the unfavourable
    deal for the UK}, in expectation: $24.7<28.6$ (computed in \S6).
\end{enumerate}

The complete ranking of the UK's sure outcomes is
$(N,\dF)\succ(N,\Nt)\succ(P,\dF)\succ(N,\dU)\succ(P,\text{hard})\succ(P,\Nt)\succ(P,\dU)\succ(N,\text{hard})$,
that is, $-19>-23>-24>-27>-27.5>-28>-32>-35$. The EU ranks
$\dU\succ\Nt\succ\dF\succ\text{hard}$ whatever the UK did.

Every deal costs the two sides 32 jointly ($19+13=23+9=27+5$). The deal type is therefore a pure
transfer between the two sides; only the probability of a hard Brexit and the preparation outlay
change total costs (\S9).

\subsection{What the payoffs leave out}

\begin{itemize}
  \item \textbf{The political cost $r$} of walking away from an offer. A government that refuses
    to put a deal to a vote may be blamed for the hard Brexit that follows; one that loses a vote
    is less exposed. The baseline sets $r=0$, and \S7.1 shows that it matters a lot.
  \item \textbf{The EU's deterrence motive.} The EU wants low costs, but it also wants to deter
    other members from leaving. A generous deal for a leaver sets a precedent. We represent this as a
    premium $k$ added to the EU's cost of an implemented favourable deal (\S7.3).
  \item \textbf{Non-monetary values}: sovereignty, the Irish border and the peace process, the
    integrity of the Union and of the single market, careers and party politics. They enter only
    through $r$ and $k$, or through the ordering.
  \item \textbf{Risk attitudes, dynamics and uncertainty}: risk neutrality is assumed; the
    future relationship and later negotiations are absent; the costs, $p$ and $r$ are known.
\end{itemize}

%-------------------------------------------------------------------------------------------
\section{(v) The game tree}

Figure~\ref{fig:tree} is the full tree of the baseline game. Filled dots are decision nodes, open
dots are Nature's (Parliament's) moves; the payoffs are (UK, EU). Thick green edges are the
subgame-perfect choices at \emph{every} node, on and off the equilibrium path (\S6).

\begin{figure}[p]
\centering
\begin{tikzpicture}[x=1.22cm,y=-0.70cm,every node/.style={font=\footnotesize},
  dec/.style={circle,fill=black,inner sep=1.7pt},
  nat/.style={circle,draw,fill=white,inner sep=1.7pt},
  e/.style={draw=gray!70,thin},
  spe/.style={draw=ans,line width=1.4pt},
  lab/.style={font=\scriptsize,fill=white,inner sep=1pt}]
  \node[dec,label={[font=\scriptsize]left:UK}] (r) at (0,8.5) {};
  \node[dec,label={[font=\scriptsize]above:EU}] (eP) at (2,4) {};
  \node[dec,label={[font=\scriptsize]below:EU}] (eN) at (2,13) {};
  \draw[spe] (r) -- node[lab,above left]{Prepare ($P$)} (eP);
  \draw[e] (r) -- node[lab,below left]{Not ($N$)} (eN);
  \branch{0}{eP}{$\dF$}{$(-24,\,-13)$}{$(-27.5,\,-18)$}{$(-27.5-r,\,-18)$}{spe}{spe}{e}
  \branch{3}{eP}{$\Nt$}{$(-28,\,-9)$}{$(-27.5,\,-18)$}{$(-27.5-r,\,-18)$}{e}{e}{spe}
  \branch{6}{eP}{$\dU$}{$(-32,\,-5)$}{$(-27.5,\,-18)$}{$(-27.5-r,\,-18)$}{e}{e}{spe}
  \branch{9}{eN}{$\dF$}{$(-19,\,-13)$}{$(-35,\,-18)$}{$(-35-r,\,-18)$}{e}{spe}{e}
  \branch{12}{eN}{$\Nt$}{$(-23,\,-9)$}{$(-35,\,-18)$}{$(-35-r,\,-18)$}{e}{spe}{e}
  \branch{15}{eN}{$\dU$}{$(-27,\,-5)$}{$(-35,\,-18)$}{$(-35-r,\,-18)$}{spe}{spe}{e}
\end{tikzpicture}
\caption{The baseline game (the EU observes preparation). At each Parliament node the upper edge
is ``pass'' (0.8) and the lower is ``fail'' (0.2). Green: the SPE action at every node, for
$r<0.4$ and $k\le5$.}\label{fig:tree}

\bigskip
\begin{tikzpicture}[x=1.3cm,y=-0.80cm,every node/.style={font=\footnotesize},
  dec/.style={circle,fill=black,inner sep=1.7pt},
  e/.style={draw=gray!70,thin},
  lab/.style={font=\scriptsize,fill=white,inner sep=1pt}]
  \node[dec,label={[font=\scriptsize]left:UK}] (r) at (0,2.5) {};
  \node[dec] (eP) at (2.2,1) {};
  \node[dec] (eN) at (2.2,4) {};
  \draw[e] (r) -- node[lab,above left]{$P$} (eP);
  \draw[e] (r) -- node[lab,below left]{$N$} (eN);
  \draw[dashed,thick,warn] (eP) -- node[lab,right,text=warn]{EU's information set} (eN);
  \foreach \yy/\dd/\pa/\txt in {0/\dF/eP/{UK presents: $(-24.7,\,-14)$},
                               1/\Nt/eP/{UK rejects: $(-27.5,\,-18)$},
                               2/\dU/eP/{UK rejects: $(-27.5,\,-18)$},
                               3/\dF/eN/{UK presents: $(-22.2,\,-14)$},
                               4/\Nt/eN/{UK presents: $(-25.4,\,-10.8)$},
                               5/\dU/eN/{UK presents: $(-28.6,\,-7.6)$}} {
    \draw[e] (\pa) -- node[lab]{$\dd$} (5.2,\yy);
    \node[right] at (5.2,\yy) {\txt};
  }
\end{tikzpicture}
\caption{The unobservable variant, reduced. The EU's two nodes are one information set. After
each offer the UK still moves at a singleton node, so each leaf shows the UK's optimal
present/reject choice (from Figure~\ref{fig:tree}, $r=0$) and the expected payoffs it
yields.}\label{fig:unobs}
\end{figure}

%-------------------------------------------------------------------------------------------
\section{(vi) Backward induction and the subgame-perfect equilibrium}

The game is finite and every information set is a singleton. Zermelo's argument (Tadelis
Proposition~8.1) carries over once each chance node is replaced by its expected payoffs, so
backward induction yields a subgame-perfect equilibrium. Every decision node starts a subgame,
so the backward-induction profile is the SPE (Tadelis \S8.2). With $r=0$ there are no ties, so the
SPE is unique. Throughout, the UK and the EU compare \emph{expected costs}.

\subsection{Step 1: the UK's present/reject decisions (six nodes)}

\textbf{After $N$.} Presenting deal $d$ costs $0.8\,c_d+0.2\times35$; rejecting costs $35+r$.
For the favourable deal:
\begin{align*}
  0.8\times19+0.2\times35 &= 15.2+7\\
                          &= 22.2 .
\end{align*}
In the same way, $\Nt$ costs $18.4+7=25.4$ and $\dU$ costs $21.6+7=28.6$. All three are below 35, and
therefore below $35+r$ for any $r\ge0$. \textbf{An unprepared UK presents every deal.}

\textbf{After $P$.} Presenting deal $d$ costs $0.8\,(c_d+5)+0.2\times27.5$; rejecting costs
$27.5+r$. For the favourable deal:
\begin{align*}
  0.8\times24+0.2\times27.5 &= 19.2+5.5\\
                            &= 24.7 .
\end{align*}
In the same way, $\Nt$ costs $22.4+5.5=27.9$ and $\dU$ costs $25.6+5.5=31.1$. With $r=0$ only $24.7<27.5$:
\textbf{a prepared UK presents $\dF$ and rejects $\Nt$ and $\dU$.}

\textbf{The general present/reject rule.} For any $p\in(0,1]$ and $r\ge0$, a prepared UK
presents $d$ when the expected cost of presenting does not exceed that of rejecting:
\begin{align}
  5+p\,c_d+(1-p)\,22.5 &\le 27.5+r \nonumber\\
  p\,c_d+22.5-22.5\,p &\le 22.5+r \nonumber\\
  p\,(c_d-22.5) &\le r \nonumber\\
  c_d &\le 22.5+\frac{r}{p}. \label{eq:rule}
\end{align}
The sunk 5 is paid in both cases, so it drops out of \eqref{eq:rule}. What moves the threshold is the
\emph{saving} $m=12.5$, which lowers the ex-post cost of no deal from 35 to 22.5. For an
unprepared UK the same steps give $c_d\le35+r/p$, which every deal satisfies.

\begin{table}[h]
\centering
\caption{Step 1: the UK's expected cost of presenting each deal versus rejecting it ($p=0.8$).}
\begin{tabular}{lcccc}
\toprule
UK at node & present $\dF$ & present $\Nt$ & present $\dU$ & reject \\
\midrule
after $N$ & \textbf{22.2} & \textbf{25.4} & \textbf{28.6} & $35+r$ \\
after $P$ & \textbf{24.7} & 27.9 & 31.1 & $\mathbf{27.5+r}$ (for $\Nt$, $\dU$) \\
\bottomrule
\end{tabular}
\end{table}

\subsection{Step 2: the EU's offer (two nodes)}

If a deal is presented, the EU's expected cost is $0.8\,e_d+0.2\times18=0.8\,e_d+3.6$:
\begin{align*}
  \dF:&\ 0.8\times13+3.6=14, & \Nt:&\ 0.8\times9+3.6=10.8, & \dU:&\ 0.8\times5+3.6=7.6 .
\end{align*}
If the UK rejects, the EU's cost is 18.
\begin{itemize}
  \item \textbf{After $N$} every deal is presented, so the EU picks the cheapest for itself:
    $\dU$ (cost 7.6). The UK's cost is 28.6.
  \item \textbf{After $P$} only $\dF$ is presented. The EU compares $\dF$ (14) with $\Nt$ or $\dU$, which are
    rejected (18), and offers $\dF$. The UK's cost is 24.7.
\end{itemize}

\subsection{Step 3: the UK's preparation decision}

Preparing leads to $(-24.7,\,-14)$ and not preparing to $(-28.6,\,-7.6)$. Since $24.7<28.6$, the UK
\textbf{prepares}. Its gain is $28.6-24.7=3.9$, and the EU's loss is $14-7.6=6.4$.

\answer{The subgame-perfect equilibrium (baseline, $r=0$, $k=0$) is the strategy profile
\begin{align*}
  s^*_{UK} &= \bigl(P;\ \text{after }P\text{: present }\dF,\ \text{reject }\Nt,\ \text{reject }\dU;\
              \text{after }N\text{: present }\dF,\Nt,\dU\bigr),\\
  s^*_{EU} &= \bigl(\dF\text{ after }P;\ \dU\text{ after }N\bigr),
\end{align*}
with Parliament passing a presented deal with probability 0.8. On the equilibrium path the UK prepares,
the EU offers $\dF$, and the UK presents it. With probability 0.8 the vote passes, giving
$(-24,\,-13)$. With probability 0.2 it fails, giving a prepared hard Brexit, $(-27.5,\,-18)$.
Expected payoffs are $(-24.7,\,-14)$. The actions after $N$ are off the path, yet they are part
of the equilibrium: they tell the UK what not preparing would cost it.}

\subsection{A Nash equilibrium that is not subgame perfect}

Consider the profile in which the UK does \emph{not} prepare but says it would reject anything but
the favourable deal, and the EU believes it:
\begin{align*}
  s'_{UK} &= \bigl(N;\ \text{after }N\text{: present }\dF,\ \text{reject }\Nt,\dU;\
              \text{after }P\text{: present }\dF,\ \text{reject }\Nt,\dU\bigr),\\
  s'_{EU} &= \bigl(\dF\text{ after }P;\ \dF\text{ after }N\bigr).
\end{align*}
The outcome is $N$, then $\dF$, presented, with expected payoffs $(-22.2,\,-14)$. This is better for the
UK than the SPE.

\textbf{It is a Nash equilibrium.} Given $s'_{EU}$, preparing would cost the UK 24.7 instead of
22.2, and presenting $\dF$ is optimal at the node that is reached. Given $s'_{UK}$, offering $\Nt$ or $\dU$
after $N$ would be rejected and cost the EU 18 instead of 14.

\textbf{It is not subgame perfect.} In the subgame after $(N,\dU)$, rejecting costs the UK 35
while presenting costs 28.6. The threat to reject is never tested on the path, so Nash
equilibrium does not discipline it; subgame perfection does (Tadelis \S8.2,
\texttt{Notas/L03/L03-05-subgame-perfection.md}). Without preparation the threat is cheap talk.

The UK's \emph{price of credibility} is the gap between this non-credible ideal and the SPE:
$24.7-22.2=2.5=5-0.2\times12.5$. That is the outlay minus its expected saving. The script
enumerates all $128\times9$ pure-strategy profiles: there are 88 pure Nash equilibria, with outcomes
$(-22.2,-14)$, $(-24.7,-14)$, $(-25.4,-10.8)$ and $(-27.5,-18)$, and exactly \textbf{one} of them is
subgame perfect. As in the ultimatum game (Tadelis Proposition~11.1), Nash equilibrium
``has no bite'' here. Non-credible threats support almost anything.

%-------------------------------------------------------------------------------------------
\section{Sensitivity}

\subsection{\texorpdfstring{The political cost of walking away, $r$}{The political cost of walking away, r}}

By \eqref{eq:rule} with $p=0.8$, a prepared UK presents $\Nt$ iff $23\le22.5+r/0.8$, that is,
$r\ge0.8\times0.5=0.4$. It presents $\dU$ iff $27\le22.5+r/0.8$, that is, $r\ge0.8\times4.5=3.6$.
The after-$N$ rule never changes. Three regimes follow (ties at the thresholds are broken
toward presenting):

\begin{center}
\begin{tikzpicture}[x=1cm]
  \draw[->,thick] (0,0) -- (13,0) node[right]{$r$};
  \foreach \x/\lb in {0/0,3/0.4,8.5/3.6} {\draw (\x,0.12) -- (\x,-0.12) node[below]{$\lb$};}
  \node[align=center,font=\footnotesize] at (1.5,0.75) {$(P,\dF)$\\UK $-24.7$, EU $-14$};
  \node[align=center,font=\footnotesize] at (5.75,0.75) {$(P,\Nt)$\\UK $-27.9$, EU $-10.8$};
  \node[align=center,font=\footnotesize] at (10.8,0.75) {$(N,\dU)$\\UK $-28.6$, EU $-7.6$};
\end{tikzpicture}

\footnotesize (not to scale)
\end{center}

\begin{itemize}
  \item \textbf{$r<0.4$}: the baseline. A prepared UK presents only $\dF$; the EU offers $\dF$.
  \item \textbf{$0.4<r<3.6$}: a prepared UK presents $\dF$ and $\Nt$, so the EU offers $\Nt$ (10.8
    beats 14). Preparing gives the UK $5+0.8\times23+0.2\times22.5=27.9$, which still beats 28.6, so the
    UK prepares. Preparation buys it only the neutral deal, a gain of $0.7$.
  \item \textbf{$r>3.6$}: even a prepared UK would present $\dU$, so the EU offers $\dU$ in both
    subgames. Preparing would cost $5+0.8\times27+0.2\times22.5=31.1>28.6$, so the UK does not
    prepare. The outcome is $(N,\dU)$.
\end{itemize}

The comparative static is the opposite of what one might expect. A higher political cost of walking
away is a cost of \emph{carrying out} the threat, so it weakens the threat. The UK wants a
low $r$, for instance a public that would not blame it for no deal. A \emph{negative} $r$, a political
reward for walking away, would make the threat credible even without preparation: an unprepared
UK rejects $\dU$ iff $28.6>35+r$, that is, iff $r<-6.4$.

\subsection{\texorpdfstring{A general pass probability $p$}{A general pass probability p}}

With $r=0$, rule \eqref{eq:rule} gives ``present $\dF$ only'' after $P$ for every $p>0$. After $N$ the
UK presents every deal and the EU offers $\dU$. The UK's gain from preparing is
\begin{align*}
  \bigl[27p+35(1-p)\bigr]-\bigl[5+19p+22.5(1-p)\bigr]
    &= (27-19-5)\,p+(35-22.5-5)(1-p)\\
    &= 3p+7.5(1-p)\\
    &= 7.5-4.5p,
\end{align*}
which is positive for every $p\in[0,1]$ and equals 3.9 at $p=0.8$. The two terms read
naturally. If the deal passes, preparation buys a better deal ($27-19=8$) for 5. If it fails,
preparation is insurance ($35-22.5=12.5$) for the same 5.

For general $p$ the thresholds in $r$ become $r\ge0.5\,p$ (present $\Nt$) and $r\ge4.5\,p$
(present $\dU$). In the middle regime the UK prepares iff $27.5+0.5p<35-8p$, that is,
$p<15/17\approx0.88$. In the high regime preparation is pure insurance, and the UK
prepares iff $27.5+4.5p<35-8p$, that is, $p<0.6$. A Parliament likely to pass deals ($p$ near 1)
makes insurance worthless and leaves only the strategic motive.

\subsection{\texorpdfstring{The EU's deterrence premium $k$}{The EU's deterrence premium k}}

Suppose an implemented favourable deal costs the EU $13+k$, where $k$ is the reputational cost of
rewarding a leaver in the eyes of other members. After $P$, the EU offers $\dF$ iff
\begin{align*}
  0.8\,(13+k)+0.2\times18 &\le 18\\
  10.4+0.8k+3.6 &\le 18\\
  0.8k &\le 4\\
  k &\le 5 .
\end{align*}
For $k>5$ the EU offers $\Nt$ or $\dU$ to a prepared UK, both are rejected, and the game ends in a
hard Brexit: $(-27.5,\,-18)$. The UK still prepares ($27.5<28.6$), so the outcome is no deal by
design, with joint cost 45.5. (If the premium were paid on \emph{offering} $\dF$ rather than on
implementing it, the condition would be $14+k\le18$, that is, $k\le4$.) After $N$ the EU offers
$\dU$ whatever $k$ is. Deterrence concerns can thus turn a commitment that works into a breakdown: a
tough EU and a credible UK are a bad combination.

\subsection{Pass probabilities that depend on the deal}

Let $p_d$ be the probability that Parliament passes deal $d$. Relative to a hard Brexit, a presented deal
saves the EU $p_d\,(18-e_d)$: $5p_{\dF}$, $9p_{\Nt}$ and $13p_{\dU}$. An unprepared UK still presents
everything, so the EU offers $\dU$ to it only if $13p_{\dU}\ge\max\{9p_{\Nt},\,5p_{\dF}\}$. Take
$(p_{\dF},p_{\Nt},p_{\dU})=(0.9,\,0.8,\,0.5)$. The savings are $4.5$, $7.2$ and $6.5$, so the EU offers
$\Nt$ after $N$ and the UK's cost is 25.4. After $P$, with $r=0$, rule \eqref{eq:rule} (with $p_d$) still
presents only $\dF$, so the EU offers $\dF$: the EU's cost is $0.9\times13+0.1\times18=13.5$ and the UK's is
$5+0.9\times19+0.1\times22.5=24.35$. The UK still prepares, but its gain falls to $1.05$. A
Parliament that is hard on bad deals does part of the work of preparation.

%-------------------------------------------------------------------------------------------
\section{The unobservable variant: commitment must be seen}

If the EU cannot see preparation, the only proper subgames start at the UK's present/reject
nodes. Those decisions are as in Step~1. Substituting them gives the reduced normal form
(rows: UK; columns: EU; payoffs (UK, EU)):

\begin{center}
\begin{tabular}{lccc}
\toprule
 & $\dF$ & $\Nt$ & $\dU$ \\
\midrule
$P$ & $(-24.7,\ \mathbf{-14})$ & $(\mathbf{-27.5},\ -18)$ & $(\mathbf{-27.5},\ -18)$ \\
$N$ & $(\mathbf{-22.2},\ -14)$ & $(\mathbf{-25.4},\ -10.8)$ & $(-28.6,\ \mathbf{-7.6})$ \\
\bottomrule
\end{tabular}
\end{center}
Bold marks best responses: the EU's within a row, the UK's within a column (for $\Nt$, the UK's best
response is $N$; after $P$, $\Nt$ and $\dU$ tie for the EU at $-18$, but $\dF$ beats both).

\textbf{No pure Nash equilibrium.} The best responses cycle. Against $P$ the EU offers $\dF$;
against $\dF$ the UK does not prepare; against $N$ the EU offers $\dU$; against $\dU$ the UK prepares.
This is the structure of Matching Pennies (\texttt{Notas/L02/L02-07-mixed-strategies.md}) and of
the Lecture~3 inspection mini case ``Work or shirk?''. Preparation only pays if the EU
expects it and concedes; but if the EU expects it, the UK would rather save the 5.

\textbf{The mixed equilibrium.} Let $\sigma=\Pr(P)$ and $\tau=\Pr(\dF)$. $\Nt$ is never used:
it costs the EU at least as much as $\dU$ in each row, and strictly more against $N$.
The EU is indifferent between $\dF$ and $\dU$ when
\begin{align*}
  -14 &= \sigma\,(-18)+(1-\sigma)(-7.6)\\
  14-7.6 &= (18-7.6)\,\sigma\\
  \sigma &= \frac{6.4}{10.4} = \frac{8}{13}.
\end{align*}
The UK is indifferent between $P$ and $N$ when
\begin{align*}
  -24.7\,\tau-27.5\,(1-\tau) &= -22.2\,\tau-28.6\,(1-\tau)\\
  28.6-27.5 &= (28.6-27.5+24.7-22.2)\,\tau\\
  \tau &= \frac{1.1}{3.6} = \frac{11}{36}.
\end{align*}
At $\sigma=8/13$, $\Nt$ would cost the EU $\tfrac{8}{13}\,18+\tfrac{5}{13}\,10.8=198/13\approx15.2>14$, so it
stays unused. Equilibrium payoffs:
\begin{align*}
  u_{UK} &= -27.5+2.8\,\tau = -27.5+\tfrac{2.8\times11}{36} = -\tfrac{1199}{45}\approx-26.64,
  & u_{EU} &= -14 .
\end{align*}

\textbf{Lessons.} (1)~Hidden commitment is no commitment. With observability the UK's cost is
24.7; without it, 26.64. Being seen to prepare is worth $35/18\approx1.94$ to the UK. The EU's
payoff, $-14$, is the same in both games. (2)~Breakdown becomes likely. A prepared UK meets the
unfavourable offer and rejects it with probability $\sigma(1-\tau)=50/117\approx0.43$. The
overall probability of a hard Brexit rises from 0.2 to $317/585\approx0.54$. (3)~The UK
therefore wants its preparation to be verifiable, and the EU, which gets $-14$ either way, has no
reason to block that. Lecture~3's point that a commitment must be observed and irreversible
(\texttt{Notas/L03/L03-06-credibility-entry-commitment.md}~\S3) is exactly this.

%-------------------------------------------------------------------------------------------
\section{Efficiency: preparation as rent-shifting}

\begin{center}
\begin{tabular}{lccc}
\toprule
Outcome & UK cost & EU cost & Joint cost \\
\midrule
$(N,\dU)$: no preparation ($r>3.6$, or the benchmark) & 28.6 & 7.6 & 36.2 \\
$(P,\dF)$: SPE, $r<0.4$ & 24.7 & 14 & 38.7 \\
$(P,\Nt)$: SPE, $0.4<r<3.6$ & 27.9 & 10.8 & 38.7 \\
Unobservable preparation, mixed equilibrium & 26.64 & 14 & 40.64 \\
$(P,\text{hard})$: $k>5$ & 27.5 & 18 & 45.5 \\
\bottomrule
\end{tabular}
\end{center}

Every deal costs 32 jointly, and a hard Brexit costs 53 jointly (40.5 if the UK prepared, before the
outlay). So
\begin{align*}
  (N,\dU):&\ 0.8\times32+0.2\times53=25.6+10.6=36.2,\\
  (P,\dF):&\ 5+0.8\times32+0.2\times40.5=5+25.6+8.1=38.7 .
\end{align*}
The difference, $2.5=5-0.2\times12.5$, is the outlay minus its expected saving. Preparation is
socially worthwhile only if $(1-p)\,12.5\ge5$, that is, $p\le0.6$. At $p=0.8$ it is a
\textbf{rent-shifting investment with a deadweight loss}: it moves 6.4 from the EU to the UK and
destroys 2.5 along the way ($6.4-3.9=2.5$).

This is the logic of strategic investment in industrial organization. An incumbent sinks capacity,
or posts a bond, not because the investment is productive but because it changes what the rival
does next (Cabral \S7.2, \texttt{Notas/L03/L03-06-credibility-entry-commitment.md}; the bond there
is worth 30 to the incumbent). There is one difference. Cabral's bond costs nothing on the
equilibrium path, whereas here the 5 is spent even when the deal passes, which is where the
deadweight loss comes from.

\textbf{A bargain the parties cannot strike.} Suppose a deal on the same frontier ($c+e=32$) gave the
UK $c^*=(24.7-0.2\times35)/0.8=22.125$. An unprepared UK would then be exactly as well off as in the
SPE, and the EU would pay $0.8\times9.875+3.6=11.5<14$. Both would gain from the EU promising
``a generous deal if you do not prepare''. But the promise is not credible: once the UK is
unprepared, the EU's best reply is $\dU$. The inefficiency comes from a commitment problem on the EU's
side. Moreover, no deal in the case's menu achieves it: $\Nt$ leaves an unprepared UK at 25.4.

%-------------------------------------------------------------------------------------------
\section{In-class decision guide}

\subsection*{Playing the UK}
\begin{enumerate}
  \item \textbf{Prepare, and make sure the EU sees it.} Visible preparation is worth 3.9 at
    $p=0.8$ (and $3p+7.5(1-p)>0$ for any $p$). Hidden preparation is worth far less (\S8).
  \item \textbf{If prepared: present only the favourable deal}, and reject the neutral and
    unfavourable ones, as long as your political cost of walking away is small ($r<0.4$).
    If $0.4<r<3.6$, present $\dF$ or $\Nt$ and reject $\dU$. If $r>3.6$, do not prepare at all.
  \item \textbf{If unprepared: present whatever is offered.} Every deal beats no deal. Saying
    otherwise is cheap talk, and a well-briefed EU will ignore it.
  \item \textbf{Expect a breakdown if the EU's deterrence concern is large} ($k>5$): you still
    prefer to have prepared ($27.5<28.6$).
\end{enumerate}

\subsection*{Playing the EU}
\begin{enumerate}
  \item \textbf{If you see preparation, offer the favourable deal}, unless you believe the UK's
    political cost of walking away is high. If you are unsure, offer $\Nt$ instead of $\dF$ only if
    $\Pr(r\ge0.4)>5/9$, and $\dU$ instead of $\dF$ only if $\Pr(r\ge3.6)>5/13$ (from
    $q\times10.8+(1-q)\times18=14$ and $q\times7.6+(1-q)\times18=14$).
  \item \textbf{If you see no preparation, offer the unfavourable deal.} It will be presented.
  \item \textbf{If you cannot tell}, the equilibrium has you offering $\dF$ with probability
    $11/36$ and $\dU$ otherwise. In practice, invest in finding out.
  \item \textbf{Weigh precedent.} If rewarding a leaver costs you more than 5 in deterrence, refuse
    the favourable deal, and accept that a hard Brexit follows.
\end{enumerate}

\subsection*{What a counterpart may do differently}
A UK negotiator may skip preparation because it looks like bad faith, or threaten rejection without
preparing. An EU negotiator may offer the neutral deal as a ``fair'' compromise, or call a bluff
because it believes $r$ is high. Each departure maps to a parameter above ($r$, $k$, $p_d$, beliefs)
or to the behavioural reasons in \S11. Record which one it was.

%-------------------------------------------------------------------------------------------
\section{(vii) Comparing the equilibrium with your negotiation}

The SPE is a benchmark for perfectly rational, money-maximizing players with common knowledge
of the tree. Real pairs often end elsewhere. The usual reasons:

\begin{itemize}
  \item \textbf{Fairness.} In ultimatum-game experiments, proposers typically offer 25--50\% of the
    pie and responders usually reject offers below about 20--25\% (Tadelis, Remark pp.~223--224,
    citing Camerer 2003; \texttt{Notas/L04/L04-05-ultimatum-game.md}~\S4). Here the ``pie'' is the
    surplus of a deal over an unprepared hard Brexit, $53-32=21$. The deals give the UK
    $16/21\approx76\%$ ($\dF$), $12/21\approx57\%$ ($\Nt$) and $8/21\approx38\%$ ($\dU$). By lab
    standards even $\dU$ is not an insulting offer, so a fairness-motivated rejection of $\dU$ by an
    unprepared UK would have to come from political humiliation (a negative component of $r$),
    not from the split itself. Conversely, an EU player may refuse $\dF$ as ``unfair'' to the EU.
  \item \textbf{Reputation and repeated interaction.} The EU deals with 27 members and future
    applicants, so a concession to a leaver is read by others. This is the $k$ of \S7.3, and the
    repeated-game logic of \texttt{Notas/L04/L04-02-carrot-and-stick.md}. The UK and the EU also
    negotiate again (trade, security, fisheries), which may soften both sides.
  \item \textbf{Uncertainty about $p$ and $r$.} If the EU does not know $r$, its offer is a screen
    (thresholds $5/9$ and $5/13$ in \S10); if the UK does not know how Parliament will vote, it may
    misjudge \eqref{eq:rule}. With private information, rational bargaining can end in breakdown
    or delay (\texttt{Notas/L04/L04-09-bargaining-private-information.md}).
  \item \textbf{Cheap talk versus commitment.} Statements like ``no deal is better than a bad deal''
    change nothing in the model unless backed by preparation or by a change in $r$. Pairs in which
    the EU believed an unprepared UK's threat played the non-credible Nash equilibrium of \S6.4.
  \item \textbf{Deadline extensions.} The model assumes a rejection means a hard Brexit now. If a
    rejection can instead lead to an extension and a new round, the game becomes one of
    alternating offers (\texttt{Notas/L04/L04-07-rubinstein-infinite-horizon.md}). The split then
    depends on patience, not only on outside options, and preparation matters less.
  \item \textbf{Domestic politics.} Votes, party leadership and elections make $p$, $r$ and even
    the UK's objective endogenous (\S1 remark: two-level games).
\end{itemize}

\textbf{Template: record your negotiation against the SPE.}

\begin{center}
\small
\begin{tabularx}{\textwidth}{>{\raggedright\arraybackslash}p{3.3cm}>{\raggedright\arraybackslash}p{4.3cm}X}
\toprule
Decision & SPE prediction (baseline) & What we did, and why (if different) \\
\midrule
UK: prepare? & Prepare, visibly & \\[14pt]
EU: what did it observe? & Preparation & \\[14pt]
EU: offer & $\dF$ (after $P$); $\dU$ would follow $N$ & \\[14pt]
UK: present or reject & Present $\dF$ & \\[14pt]
Parliament (dice / draw) & Pass with prob.\ 0.8 & \\[14pt]
Final outcome, payoffs & $(-24,-13)$ or $(-27.5,-18)$; exp.\ $(-24.7,-14)$ & \\[14pt]
Beliefs about $r$, $k$, $p$ & $r=0$, $k=0$, $p=0.8$ & \\[14pt]
Main reason for any gap & --- & \\[14pt]
\bottomrule
\end{tabularx}
\end{center}

%-------------------------------------------------------------------------------------------
\section{Historical note (general knowledge, not from the case)}

The facts in this section are general knowledge and are not taken from the case. Boris Johnson
became Prime Minister in July 2019. In October 2019 his government and the EU agreed a revised
Withdrawal Agreement, which replaced the Irish ``backstop'' with the Protocol on Ireland/Northern
Ireland. The Article~50 deadline had been extended three times during 2019. After the December 2019
general election, Parliament approved the Withdrawal Agreement, and the UK left the EU on
31~January~2020, entering a transition period that ended on 31~December~2020. Michel Barnier led the EU's
side in both negotiations. The \textbf{EU--UK Trade and Cooperation Agreement} was agreed on
24~December~2020, applied provisionally from 1~January~2021, and entered into force on 1~May~2021.
It provides for zero tariffs and zero quotas on goods that meet its rules of origin, and covers
much less of trade in services. Which of the model's deal types, if any, these agreements
correspond to is a question for class discussion, not something the model settles.

%-------------------------------------------------------------------------------------------
\section*{References}
\begin{itemize}[leftmargin=1.4em]\small\raggedright
  \item Tadelis, S. (2013). \emph{Game Theory: An Introduction}. Princeton University Press.
    Ch.~7, ``Preliminaries'' of dynamic games (pp.~129--150): extensive form, Nature, information
    sets, strategies. Ch.~8, ``Credibility and Sequential Rationality'' (pp.~151--174): backward
    induction, Proposition~8.1, subgame perfection. Ch.~6 (pp.~101--114): mixed strategies. \S11.1
    (pp.~220--224): the ultimatum game and its experimental Remark.
  \item Cabral, L. (2017). \emph{Introduction to Industrial Organization}, 2nd ed. MIT Press.
    \S7.2 (pp.~167--171): credible commitment and entry deterrence.
  \item Schelling, T. C. (1960). \emph{The Strategy of Conflict}. Harvard University Press.
  \item Putnam, R. D. (1988). Diplomacy and domestic politics: the logic of two-level games.
    \emph{International Organization} 42(3), 427--460.
  \item Course notes (Lectures 2--4):
    \begin{itemize}[leftmargin=1.2em,itemsep=0pt,topsep=1pt]
      \item \texttt{Notas/L03/}: \texttt{L03-01-extensive-form.md}, \texttt{L03-02-strategies.md},
        \texttt{L03-04-sequential-rationality-backward-induction.md},
        \texttt{L03-05-subgame-perfection.md}, \texttt{L03-06-credibility-entry-commitment.md};
      \item \texttt{Notas/L02/}: \texttt{L02-07-mixed-strategies.md};
      \item \texttt{Notas/L04/}: \texttt{L04-02-carrot-and-stick.md},
        \texttt{L04-05-ultimatum-game.md}, \texttt{L04-07-rubinstein-infinite-horizon.md},
        \texttt{L04-08-legislative-bargaining.md}, \texttt{L04-09-bargaining-private-information.md};
      \item Lecture~3 mini case ``Work or shirk?'': \texttt{Aula/Lecture03.pdf}, slides 41--42.
    \end{itemize}
  \item Cases: ``Boris's Brexit Dilemma (A)'' and ``(B)'', Darden UV9621 and UV9624 (figures only,
    not reproduced).
  \item Check script: \texttt{Cases/Brexit/brexit-check.py} (exact fractions; backward induction
    and normal-form enumeration on the tree).
\end{itemize}

\end{document}
